\documentclass[11pt]{article}
\usepackage{fontspec}
\usepackage[french,english]{babel}
\usepackage{amsmath}
\usepackage{amssymb}
\usepackage{amsthm}
\usepackage{mathrsfs}
\usepackage{tikz-cd}
\usepackage{enumitem}
\usetikzlibrary{babel}

\newcommand{\printedpage}[1]{\marginpar{\footnotesize p.~#1}}
\newcommand{\msp}[1]{\marginpar{\footnotesize [#1]}}
\newcommand{\Ob}{\mathsf{Ob}}
\newcommand{\id}{\mathsf{id}}
\newcommand{\DD}{\mathbb{D}}

\begin{document}

\printedpage{29}
% head: A. Grothendieck : Dérivateurs, Ch. I, Édition des Universités de Montpellier II et Paris VII
% foot: Version 8 juillet 2001

\noindent\msp{page 51}
base change for
\[
\begin{tikzcd}
Z' \arrow[r, "p"] \arrow[d, "f'"'] & Z \arrow[d, "f"] \\
S \arrow[r, "\id"] & S
\end{tikzcd}
\qquad .
\]

\noindent\textsc{Remark} : We recover the fact that $i^*i_*(F_Y) \overset{\sim}{\to} F_Y$ is \texttt{[an]} isomorphism, for $i$ a closed immersion. Moreover, with the notation of \textbf{Der 4} (g), $j^*(i_*(F_Y))$ is computed by the base change theorem (corollary 2 \texttt{[rather 4]} above), which applies in every commutative situation
\[
\begin{tikzcd}[row sep=large, column sep=large]
X' \arrow[r, "q"] \arrow[d, "f'"'] \arrow[dr, phantom, "\text{\scriptsize cart}"] & X \arrow[d, "f"] \\
Y' \arrow[r, "p"] & Y
\end{tikzcd}
\]
when $p$ is an open immersion (and $\gamma = \id$) (then $X'/y' \to X/p(y')$ is even an isomorphism!), and we find, returning to the notation of \textbf{Der 4} (g),
\[
\begin{tikzcd}
\emptyset \arrow[r, "q"] \arrow[d, "f'"'] & Y \arrow[d, hook, "i"] \\
U \arrow[r, hook, "j"'] & X
\end{tikzcd}
\]
\textit{i.e.} $j^*(i_*(F_Y)) \simeq f'_*(\varphi)$, where $\varphi = q^*(F_Y)$ is the unique element of the category (equivalent to a one-point category, by \textbf{Der 1}) $\DD(\emptyset)$. On the other hand, $f'_*$ takes final object to final object, whence $j^*i_*(F_Y)$ is a final object.

Conversely, let $F$ be in $\DD(X)$, such that $j^*(F) = e_{\DD(U)}$. Consider
\[
F \xrightarrow{u} i_*i^*F \qquad\quad \text{(adjunction) ;}
\]
I say that this is an isomorphism. Indeed, by the axiom \msp{page 52} of localization \textbf{Der 2}, it suffices to prove that $i^*(u)$ and $j^*(u)$ are \texttt{[some]} isomorphisms. For $i^*(u)$ this follows from the fact that $i^*$ is a localization functor, and for $j^*$ from the fact that $j^*(F)$ and $j^*(i_*i^*(F))$ are final objects in $\DD(U)$.

\printedpage{30}
% head: A. Grothendieck : Dérivateurs, Ch. I, Édition des Universités de Montpellier II et Paris VII
% foot: Version 8 juillet 2001

\noindent\msp{page 53}

\subsection*{Homotopically proper, homotopically smooth morphisms.}

\noindent Let $f : X \to Y$ be \texttt{[a]} morphism of categories, consider $y_0 \in \Ob\, Y$, and the canonical functor
\[
X_{y_0} \xrightarrow{\varphi} X/y_0 \qquad\qquad\qquad \varphi(x) = (x, \id_{f(x)}) \tag*{(1)}
\]
($X_{y_0}$ fiber of $X$ at $y_0$). Let $(x_0, u_0) \in \Ob\, X/y_0$ hence $x_0 \in \Ob\, X$, $u_0 : f(x_0) \to y_0$. Consider the category coinduced by $\varphi$ over $(x_0, u_0)$
\[
(x_0, u_0)\backslash(X_{y_0}) \quad ,
\]
to test whether $\varphi$ is a costrict hot-equivalence \textit{i.e.} homologically \texttt{[rather co\-ho\-mo\-lo\-gi\-cal\-ly]} cofinal. We are going to interpret it via the functor induced by $f$
\[
x_0\backslash X \xrightarrow{\psi} f(x_0)\backslash Y
\]
going from the category of specializations of $x_0$ in $X$ to that of specializations of $f(x_0)$ in $Y$. We then find a canonical isomorphism between the category coinduced $(x_0, u_0)\backslash X_{y_0}$ by $X_{y_0}$ over the object $(x_0, u_0)$ of $X/y_0$ and the fiber category $(x_0\backslash X)_{u_0}$ of $x_0\backslash X$ over $f(x_0)\backslash Y$ at $u_0$. Thus

\noindent\textbf{Proposition.} \textit{For the functors $\varphi : X_{y_0} \to X/y_0$ ($y_0 \in \Ob\, Y$) to be costrict hot-equivalences, it is necessary and sufficient that the fibers of the functors $x_0\backslash X \to f(x_0)\backslash Y$ induced by $f$ be contractible.}

\noindent\msp{page 54}
We then say that $f$ is \textit{homotopically proper.} The dual notion is expressed, either by the condition that the functors induced by $f$
\[
X_{y_0} \longrightarrow y_0\backslash Y
\]
be strict hot-equivalences \textit{i.e.} that they be cohomologically \texttt{[rather ho\-mo\-lo\-gi\-cal\-ly]} cofinal, or that the fibers of the morphisms induced on the ``localized categories'' (categories of generizations)
\[
X/x_0 \longrightarrow Y/f(x_0)
\]
be contractible. (N.B. We use here that a category $Z$ is contractible if and only if $Z^\circ$ is.) We then say that $f$ is \textit{homotopically smooth.}

\noindent\textbf{Proposition.} \textit{The notion of proper (resp. smooth) functor is stable under base change, under composition. Isomorphisms, and more generally, closed immersions are proper. The image of a proper morphism is a closed subcategory.}\footnote[23]{Beware, an equivalence of categories is not proper! If $X' \hookrightarrow X$ \texttt{[is]} an inclusion of a subcategory of $X$, it is proper only if it is a closed immersion.}

\noindent N.B. The functor $X \to e$ is always proper and smooth, for the categories $x_0\backslash X$ and $X/x_0$ are contractible (having an initial object or a final object). Thus every projection functor $X \times Y \xrightarrow{pr_2} Y$ is both proper and smooth.

\printedpage{31}
% head: A. Grothendieck : Dérivateurs, Ch. I, Édition des Universités de Montpellier II et Paris VII
% foot: Version 8 juillet 2001

\noindent\msp{page 55}
\textsc{Proof}. Stability under base change :
\[
\begin{tikzcd}
X \arrow[d, "f"'] \arrow[dr, phantom, "\text{\scriptsize cart}"] & X' \arrow[l, "q"'] \arrow[d, "f'"] \\
Y & Y' \arrow[l, "p"]
\end{tikzcd}
\qquad ,
\]
whence for $x' \in \Ob\, X'$
\[
\begin{tikzcd}[row sep=large]
X/x \arrow[d] & X'/x' \arrow[l] \arrow[d] \\
Y/y & Y'/y' \arrow[l]
\end{tikzcd}
\qquad ,
\]
where $x$, $y'$, $y$ are the images of $x'$ in $X$, $Y'$, $Y$. We check that this diagram is \textit{cartesian}, hence the fibers of $X'/x' \longrightarrow Y'/y'$ are isomorphic to fibers of $X/x \longrightarrow Y/y$, hence contractible if these are.

In a case $X \xrightarrow{f} Y$ we check that for a bilocalized functor of $f$
\[
X/x \xrightarrow{\varphi=f/x} Y/y \quad ,
\]
every ``bilocalized''
\[
(X/x)/\xi \longrightarrow (Y/y)/\eta \qquad\qquad (\xi : x' \to x\,, \quad \eta : y' \to y)
\]
of $\varphi$ is again identified with a bilocalized of $f$
\[
X/x' \longrightarrow Y/y' \quad .
\]
Thus if $f$ is smooth, the bilocalized $f/x$ ($x \in \Ob\, X$) are so too. Let us note on the other hand.

\noindent\textbf{Proposition.}\footnote[24]{\textit{cf.} improvement in cor. 2, next page.} \textit{Every proper (or dually smooth) morphism with contractible fibers is a hot-equivalence.}

\noindent\msp{page 56}
This comes from the fact that for such a proper functor, if $F$ is a presheaf, or complex of abelian presheaves on $X$, the formation of the $Rf_*(F)$ commutes with passage to fibers, and likewise for $R^1f_*(F)$ if $F$ is a presheaf of groups (not even abelian).

\noindent\textbf{Corollary 1.} \textit{If $f : X \to Y$ is cohomologically proper with contractible fibers, then $Y$ is contractible if and only if $X$ is.}

Thanks to this, we find that the composite of two proper (smooth) functors is proper (smooth).

\printedpage{32}
% head: A. Grothendieck : Dérivateurs, Ch. I, Édition des Universités de Montpellier II et Paris VII
% foot: Version 8 juillet 2001

\noindent N.B. Contrary to what happens in topology for the notion of properness, if $gf$ is proper and $g$ proper, $f$ is not necessarily so, \textit{i.e.} an $S$-morphism of $S$-categories $X$, $Y$ proper over $S$ is not necessarily proper. Thus a section of $X/S$, for $X$ proper over $S$, is not necessarily proper, \textit{i.e.} is not necessarily a closed immersion. This is already false if $S = e$ (in which case $f : X \to e$ is always proper).

\noindent\textbf{Corollary 2.} \textit{A proper or smooth morphism, with contractible fibers is both a strict and a costrict hot-equivalence.}

Indeed, we apply corollary 1 to $X/y \to Y/y$ or $y\backslash X \to y\backslash Y$.

\noindent\textbf{Theorem} (g) \textit{Consider a base change situation in} \textsf{Dia}
\[
\begin{tikzcd}[row sep=large, column sep=huge]
X \arrow[d, "f"'] \arrow[dr, phantom, "\text{cart}"] & X' \arrow[l, "q"'] \arrow[d, "f'"] \\
Y & Y' \arrow[l, "p"]
\end{tikzcd}
\qquad , \qquad\qquad p^*f_*(F) \xrightarrow{c} f'_*q^*(F)
\]
\textit{($F \in \DD(X)$), $\DD$ a left exact derivator. If $f$ is (homotopically) proper or $p$ homotopically smooth, then $c$ is an isomorphism.}

\noindent\msp{page 57}
The question is whether for every $y' \in Y'$, putting $y = p(y')$ and considering the functor induced by $q$
\[
X/y \longleftarrow X'/y'
\]
the corresponding morphism
\[
\Gamma_{\DD}(X/y, F) \longrightarrow \Gamma_{\DD}(X'/y', q^*F)
\]
is \texttt{[an]} isomorphism. In the case $f$ proper, we consider
\[
\begin{tikzcd}[column sep=large]
X_y \arrow[d, "\sim\ \ \text{costrict}"'] & X_{y'} \arrow[l, "\sim"'] \arrow[d, "\sim\ \ \text{costrict}"] \\
X/y & X'/y' \arrow[l]
\end{tikzcd}
\]
whence the conclusion, since by \textbf{Der 5} (g), in
\[
\begin{tikzcd}[row sep=large, column sep=large]
\Gamma(X_y) \arrow[r] & \Gamma(X_{y'}) \\
\Gamma(X/y) \arrow[u] \arrow[r] & \Gamma(X'/y') \arrow[u]
\end{tikzcd}
\]
the vertical arrows and the upper horizontal one are \texttt{[some]} isomorphisms, hence also the lower horizontal one.

\printedpage{33}
% head: A. Grothendieck : Dérivateurs, Ch. I, Édition des Universités de Montpellier II et Paris VII
% foot: Version 8 juillet 2001

In the case $p$ smooth, let us consider on the $Y$, $Y'$ side the base change square % [sic] coté -> côté, missing circumflex (cf. "Du côté" below)
\[
\begin{tikzcd}[row sep=large, column sep=large]
Y & Y' \arrow[l] \\
Y/y \arrow[u] & Y'/y' \arrow[l] \arrow[u]
\end{tikzcd}
\]

\noindent\msp{page 58}
whence the adjoining diagram
\[
\begin{tikzcd}[column sep=small]
X \arrow[dd] & & X' \arrow[ll] \arrow[dd] & \\
& X/y \arrow[ul] & & X'/y' \arrow[ll, crossing over, "\kappa"'] \arrow[ul] \arrow[dd] \\
Y & & Y' \arrow[ll] & \\
& Y/y \arrow[ul] \arrow[from=uu, crossing over] & & Y'/y' \arrow[ll, "\pi"] \arrow[ul]
\end{tikzcd}
\qquad ,
\]
where the lateral (vertical) squares are cartesian. We note that $\pi$ is smooth and has contractible fibers (for $p$ is smooth) hence $\kappa$ also. Thus $\kappa$ is a costrict hot-equivalence, hence $\Gamma(X/y) \overset{\sim}{\longrightarrow} \Gamma(X'/y')$, we win again.

\noindent N.B. Unless I am mistaken, we have a converse of the theorem, \texttt{[to]} wit : If $f : X \to Y$ is such that the base change theorem for $f_*(F)$ holds for \textit{every} derivator (it probably suffices to take \textsf{HOTAB} and \textsf{HOT}) and for every $Y' \to Y$, then $f$ is homotopically proper. On the other hand, if $g : Y' \to Y$ is such that the base change theorem for $g$ holds for every $X \to Y$ and every $F \in \DD(X)$ ($\DD$ any right exact derivator), then $g$ is homotopically smooth. % [sic] pout tout -> pour tout, misprint

\noindent On the homological side the smooth functors $f$ are exactly those for which $f_!$ commutes with every base change. The proper functors $g$ are exactly those for which base change by $g$ commutes with $f_!$, for every $f : X \to Y$ and every $\DD$ a left exact derivator.

\noindent\msp{page 59}
Thus, the proper functors $f : X \to Y$ are characterized in \textit{four} remarkable ways
\begin{itemize}[label={}, leftmargin=0pt, itemindent=2\parindent, labelsep=0.3em]
\item[\textit{a})] $f_*$ commutes with every base change ($\forall$ left exact derivator $\DD$);
\item[\textit{b})] the base change $f^*$ commutes with every $g_!$ ($g : Y' \to Y$, $\DD$ a right exact derivator);
\item[\textit{c})] $\forall y \in \Ob\, Y$, $X_y \to X/y$ is a costrict hot-equivalence \textit{i.e.} is a homologically \texttt{[rather co\-ho\-mo\-lo\-gi\-cal\-ly]} cofinal functor;
\item[\textit{d})] $\forall x \in \Ob\, X$ the functor induced on the colocalized
\[
x\backslash X \longrightarrow f(x)\backslash Y
\]
\end{itemize}

\end{document}
